
Let us now prove \hyperref[thm:sga3-V-4.1-iii]{Theorem~4.1(iii)}. Recall once more that
\(\underline P\) denotes the set underlying a scheme \(P\), and that
\(\underline d:\underline P\to\underline Q\) is the map induced by a
morphism \(d:P\to Q\). \hyperref[proof:sga3-V-4.1-b]{Part~b)} can then be restated by saying that the
map
\[
  \underline{d_0}\boxtimes\underline{d_1}:
  \underline{X_1}\longrightarrow
  \underline{X_0}\times_{\underline Y}\underline{X_0}
\]
with components \(\underline{d_0}\) and \(\underline{d_1}\) is
surjective. This map factors as
\[
  \underline{X_1}
  \xrightarrow{\ \underline{d_0}\boxtimes\underline{d_1}\ }
  \underline{X_0}\times_{\underline Y}\underline{X_0}
  \xrightarrow{\ q\ }
  \underline{X_0\times_YX_0},
\]
where \(q\) is the canonical map. The image of
\(\underline{d_0}\boxtimes\underline{d_1}\) therefore contains every
point \(v\in\underline{X_0\times_YX_0}\) such that
\[
  \{v\}=q^{-1}(q(v)).
\]
This condition is satisfied in particular when \(v\) is
\emph{rational over \(Y\)}, that is, when the residue field
\(\kappa(v)\) is identified with the residue field \(\kappa(w)\) of
the image \(w\) of \(v\) in \(Y\).%

If \(v\in\underline{X_0\times_YX_0}\) is not rational over \(Y\),
again let \(w\)
be its image in \(Y\). By
Lemma~\hyperref[lem:sga3-V-4.1.2]{4.1.2}, there exist a local ring
\(C\) with radical \(\mathfrak m\) and a local, flat homomorphism
\[
  f:\mathcal O_{Y,w}\longrightarrow C
\]
such that \(C/\mathfrak m\) is isomorphic to \(\kappa(v)\) as a
\(\kappa(w)\)-algebra. Put \(Y'=\operatorname{Spec}C\), and let
\(\pi:Y'\to Y\) be the morphism induced by \(f\). The canonical
projection
\[
  (X_0\times_YX_0)\times_YY'\longrightarrow X_0\times_YX_0
\]
sends some point \(v'\), rational over \(Y'\), to \(v\). Since
\[
  (X_0\times_YX_0)\times_YY'
  \simeq
  (X_0\times_YY')\times_{Y'}(X_0\times_YY'),
\]
and since the hypotheses of \hyperref[thm:sga3-V-4.1]{Theorem~4.1} and the preceding results,
in particular \hyperref[proof:sga3-V-4.1-b]{part~b)}, remain valid after the base change
\(\pi:Y'\to Y\), the point \(v'\) is the image of an element
\[
  u'\in\underline{X_1\times_YY'}
\]
under the base change of \(d_0\boxtimes d_1\). If \(u\) is the image
of \(u'\) in \(X_1\), then
\[
  v=(d_0\boxtimes d_1)(u),
\]
as required.

\medskip
\phantomsection\label{proof:sga3-V-4.1-f}%
\noindent\textbf{f) Proof of \hyperref[thm:sga3-V-4.1-iv]{(iv)}.}%
\footnote{\label{note:sga3-V-24}Editor's note: \hyperref[lem:sga3-V-4.1.3]{the following lemma} has been added from
\hyperref[bib:sga3-V-DG70]{[DG70]}, I, \S5, Proposition~1.5; see also the proof of EGA
\(\mathrm{IV}_3\), 8.11.5. It was used implicitly in the original and
explicitly in \hyperref[bib:sga3-V-DG70]{[DG70]}, III, \S2, 4.6. It is also useful in
\hyperref[thm:sga3-V-7.1]{Theorem~7.1 below}.}

\medskip
\phantomsection\label{lem:sga3-V-4.1.3}%
\noindent\textbf{Lemma 4.1.3.}---
\textit{If a monomorphism of schemes \(f:T\to Z\) is finite, it is a
closed immersion.}

Covering \(Z\) by affine open subsets \(Z_i\), and replacing \(f\) by
the induced morphisms \(f^{-1}(Z_i)\to Z_i\), reduces us---since a
finite morphism is affine---to the case
\[
  Z=\operatorname{Spec}B,\qquad T=\operatorname{Spec}A.
\]
Because \(f\) is a monomorphism, the diagonal morphism
\[
  T\longrightarrow T\times_ZT
\]
is an isomorphism (EGA~I, 5.3.8); equivalently, multiplication gives
an isomorphism
\[
  A\otimes_BA\xrightarrow{\ \sim\ }A.
\]
Consequently, for every maximal ideal \(\mathfrak m\) of \(B\), there
is an isomorphism
\[
  (A/\mathfrak mA)\otimes_k(A/\mathfrak mA)
  \xrightarrow{\ \sim\ }A/\mathfrak mA,
  \qquad k=B/\mathfrak m.
\]
Since \(A\) is finite over \(B\), the \(k\)-vector space
\(A/\mathfrak mA\) has finite dimension \(d\), and this isomorphism
implies \(d=0\) or \(d=1\). Hence the map
\[
  k=B/\mathfrak m\longrightarrow A/\mathfrak mA
\]
is surjective. Nakayama's lemma, applied after localization at
\(\mathfrak m\), shows that \(B_{\mathfrak m}\to A_{\mathfrak m}\)
is surjective. It follows that \(B\to A\) is surjective: its cokernel
\(C\) satisfies \(C_{\mathfrak m}=0\) for every maximal ideal
\(\mathfrak m\), and therefore vanishes. This proves \hyperref[lem:sga3-V-4.1.3]{the lemma}.
